\begin{lemma}[Minimal deciding prefixes form a front]
If $h$ is locally constant, then $F^h$ from \noderef{DecidingFrontSet} is a
front on $\mathbb N$.
\end{lemma}
\begin{proof}
Write $F=\mathsf{DecidingFrontSet}(h)$, using
\noderef{DecidingFrontSet}, and let $S^h$ denote the set of finite sets $s$
satisfying $\mathsf{DecidingPrefix}(h,s)$, as defined in
\noderef{DecidingPrefix}.  We first establish the finite-prefix fact used
throughout the proof.  For $x\in\mathsf{IncSeq}$, as defined in
\noderef{IncSeq}, and $n\in\mathbb N$, put
\[
 P_n(x)=\{x(i)\mid i<n\}.
\]
Because $x$ is an order embedding, its values are pairwise distinct and
strictly increasing.  Therefore the values of $x$ below $x(n)$ are exactly
$x(0),\ldots,x(n-1)$: the forward implication follows from strict increase,
and if $x(i)<x(n)$ with $n\leq i$, strict increase would give
$x(n)\leq x(i)$, a contradiction.  Thus $x(n)$ witnesses
\[
 \mathsf{ProperPrefixSet}(P_n(x),\operatorname{range}(x)),
\]
using the defining cutoff equation in \noderef{ProperPrefixSet}.

We also need the following agreement consequence.  If
\[
 \mathsf{ProperPrefixSet}(P_n(x),\operatorname{range}(y))
\]
for another increasing enumeration $y$, unpacking
\noderef{ProperPrefixSet} gives a cutoff $y(m)$ such that
$P_n(x)=\{y(i)\mid i<m\}$.  The first set has $n$ elements and the second
has $m$ elements, since both sequences are injective, so $m=n$.  The two
finite sets of the first $n$ values are consequently equal.  We now prove
by induction on $i<n$ that $x(i)=y(i)$.  At $i=0$, both values are the least
element of this common finite set.  If the equality has been proved below
$i$, then $x(i)$ and $y(i)$ are both the least member of the common set not
among the already equal values; strict increase makes this least-member
description unique.  This proves
\[
 \mathsf{ProperPrefixSet}(P_n(x),\operatorname{range}(y))
 \quad\Longrightarrow\quad
 \mathsf{PrefixAgree}(n,x,y)
\tag{1}
\]
The same argument applies when the common finite prefix is written simply
as $s$: if $s=P_n(x)$ and both $s$ and $s$ are proper prefixes of the ranges
of $x$ and $y$, then (1) gives $\mathsf{PrefixAgree}(n,x,y)$.

Now fix $x\in\mathsf{IncSeq}$.  By \noderef{LocallyConstant}, choose
$n$ such that every $y$ satisfying
$\mathsf{PrefixAgree}(n,x,y)$, as defined in \noderef{PrefixAgree}, has
$h(y)=h(x)$.  Set $s=P_n(x)$.  The first
paragraph shows that $s$ is a proper prefix of $\operatorname{range}(x)$.
If $y$ and $z$ both have $s$ as a proper prefix of their ranges, (1) gives
$\mathsf{PrefixAgree}(n,x,y)$ and
$\mathsf{PrefixAgree}(n,x,z)$.  Hence
\[
 h(y)=h(x)=h(z).
\]
By the definition of $\mathsf{DecidingPrefix}$ in \noderef{DecidingPrefix},
this proves that $s$ belongs to $S^h$, so every branch has a deciding finite
prefix.

We next make the minimality step explicit.  If $s$ satisfies
$\mathsf{DecidingPrefix}(h,s)$, then there is $t\in F$ with
$\mathsf{InitialSegment}(t,s)$, as defined in \noderef{InitialSegment}.
Prove this by induction on the finite
cardinality of $s$.  If $s\in F$, take $t=s$ and use the equality clause of
\noderef{InitialSegment}.  Otherwise, unfolding
\noderef{DecidingFrontSet} and using that $s$ is deciding gives a $q$ with
 $\mathsf{ProperInitialSegment}(q,s)$, as defined in
\noderef{ProperInitialSegment}, and $q$ deciding.  The proper
initial-segment description in \noderef{ProperInitialSegment} and
\noderef{InitialSegment} shows that $q$ is a proper subset of $s$, hence
$|q|<|s|$.  The induction hypothesis supplies $t\in F$ with
 $\mathsf{InitialSegment}(t,q)$.  Composing the equality-or-cutoff
descriptions in \noderef{InitialSegment} shows
$\mathsf{InitialSegment}(t,s)$: an initial segment of an initial segment
is an initial segment.  This proves the claim.

Let now $Y\subseteq\mathbb N$ be infinite.  By
\noderef{InfiniteSetEnumeration}, choose $x\in\mathsf{IncSeq}$ with
$\operatorname{range}(x)=Y$.  The preceding local-constancy argument gives
a deciding $s=P_n(x)$, and the minimality claim gives $t\in F$ with
$\mathsf{InitialSegment}(t,s)$.  We claim that
$\mathsf{ProperPrefixSet}(t,Y)$, using \noderef{ProperPrefixSet}.  If $t=s$, this is the already established
proper-prefix assertion for $s$.  Otherwise, \noderef{InitialSegment} and
the inequality $t\ne s$ give $\mathsf{ProperInitialSegment}(t,s)$, as
defined in \noderef{ProperInitialSegment}, so
$t=\{k\in s\mid k<b\}$ for some $b\in s$.  Since
$s=\{k\in Y\mid k<a\}$ for the cutoff $a=x(n)\in Y$, the same equation
gives $t=\{k\in Y\mid k<b\}$, with $b\in Y$.  This is exactly
\noderef{ProperPrefixSet}.  Thus the family $F$ is dense in the sense of
the density clause of \noderef{Front}.

For prefix-freeness, let $s,t\in F$ and assume
$\mathsf{InitialSegment}(s,t)$.  If $s\ne t$, then
$\mathsf{ProperInitialSegment}(s,t)$, while membership of $t$ in
$\mathsf{DecidingFrontSet}$, as defined in \noderef{DecidingFrontSet}, says that no proper initial segment of $t$ is a
member of $S^h$.  But $s\in F$ implies, by the first conjunct in the
definition of \noderef{DecidingFrontSet}, that $s$ is deciding.  This is a
contradiction.  Therefore $s=t$, which is precisely the prefix-free clause
of \noderef{Front}.

It remains to verify the base alternative.  First suppose that $h$ is
constant, meaning $h(x)=h(y)$ for all $x,y\in\mathsf{IncSeq}$.  For every
$x$, the empty set is a proper prefix of $\operatorname{range}(x)$: use the
least range element $x(0)$ as the cutoff in \noderef{ProperPrefixSet}.
Consequently the constant-value hypothesis proves
 $\mathsf{DecidingPrefix}(h,\emptyset)$, by \noderef{DecidingPrefix}.  No proper initial segment of the
empty set exists, because the non-equality alternative in
\noderef{InitialSegment} requires a cutoff element of the right-hand finite
set.  Hence $\emptyset\in F$ by the definition of
\noderef{DecidingFrontSet}.  Conversely, if $s\in F$ and $s\ne\emptyset$,
let $a$ be the least member of $s$.  Then
$\emptyset=\{k\in s\mid k<a\}$, so $\mathsf{InitialSegment}$, as defined
in \noderef{InitialSegment}, and $s\ne\emptyset$ give
$\mathsf{ProperInitialSegment}(\emptyset,s)$ by
\noderef{ProperInitialSegment}.
This contradicts the minimality condition defining $F$, since
$\emptyset$ is deciding.  Thus every member of $F$ is empty and
$F=\{\emptyset\}$.

Conversely, suppose $h$ is not constant.  If the empty set were deciding,
then the least-element argument just given would show that the empty set is
a proper prefix of the range of every $x\in\mathsf{IncSeq}$.  The defining
condition \noderef{DecidingPrefix} would then give $h(x)=h(y)$ for all
$x,y$, contrary to the assumption.  Therefore
$\neg\mathsf{DecidingPrefix}(h,\emptyset)$.

Fix $n\in\mathbb N$ and let
\[
 Y_n=\{m\in\mathbb N\mid n\leq m\}.
\]
This set is infinite, and \noderef{InfiniteSetEnumeration} supplies an
increasing enumeration $x_n$ with range $Y_n$.  Applying the density
construction above to $Y_n$ gives $s\in F$ with
$\mathsf{ProperPrefixSet}(s,Y_n)$, by \noderef{ProperPrefixSet}.  The set $s$ is nonempty, since an empty
member of $F$ would be deciding by the first conjunct of
\noderef{DecidingFrontSet}, contradicting the preceding paragraph.  Unpack
\noderef{ProperPrefixSet} and choose a cutoff $a\in Y_n$ with
$s=\{k\in Y_n\mid k<a\}$.  If $a=n$, then $s=\emptyset$, so $a>n$.
Hence $n\in s$, and the definition of \noderef{FrontBase} gives
$n\in\mathsf{FrontBase}(F)$.  Since $n$ was arbitrary,
$\mathbb N\subseteq\mathsf{FrontBase}(F)$; the reverse inclusion is
immediate because every member of a finite subset of $\mathbb N$ is a
natural number.  Therefore
\[
 \mathsf{FrontBase}(F)=\mathbb N.
\]

Finally, $\mathbb N=\operatorname{Set.univ}$ is infinite (the identity
enumeration has infinitely many distinct values).  We have explicitly shown
the infinite-base clause, the base alternative (either
$F=\{\emptyset\}$ or $\mathsf{FrontBase}(F)=\operatorname{Set.univ}$), the
prefix-free clause, and the density clause.  By the four fields in the
definition of \noderef{Front}, these facts assemble to
\[
 \mathsf{Front}(\mathsf{DecidingFrontSet}(h),\operatorname{Set.univ}),
\]
as required.
\end{proof}
